%قبل از شروع مسئله ترجیحاً چیزی ننویسید. اما اگر خواستید بنویسید، در قالب لیستی (فقرات و شمارش) بنویسید.
%برای اضافه کردن نمره، بعدی از شروع‌مسئله، [x] را اضافه کنید و به جای x نمره مورد نظر را قرار دهید. برای زیرمسئله‌ها هم به همین صورت است.
%\شروع‌مسئله[10]
\شروع‌مسئله [15]
%\متن‌سیاه{سوال عملی اول:}\\
%\متن‌سیاه{دراپ‌تست}\\

% صورت اصلی مسئله را در این‌جا بنویسید.
$k$ 
تا عدد $a_1,a_2,...,a_k$ و $n$ پیشنهاد داریم که هر پیشنهاد می‌تواند یکی از دو نوع زیر باشد:
\شروع{شمارش}[1. ]
\فقره عدد $b$ به $a_i$ اضافه شود. 
\فقره عدد $b$ در $a_i$ ضرب شود. 
\پایان{شمارش}
باید
$m$
تا از این پیشنهاد‌ها را انتخاب کنید و به یک ترتیب دلخواه آن‌ها را اعمال کنید، به طوری که
$\prod_{i=1}^{k}a_i$
ماکسیمم شود.\\
\متن‌سیاه{راهنمایی:} برای مرتب کردن اعداد کسری مثل 
$\frac{a}{b}$
به طوری که $a$ و $b$ هر دو عدد صحیح هستند، اگر از متغیر
float
استفاده کنیم و دو عضو را مقایسه کنیم، امکان خطا وجود دارد، برای از بین بردن خطا می‌توان از این روش استفاده کرد: هر زمان به نتیجه 
$\frac{a}{b}\leq \frac{a'}{b'}$
نیاز داشتیم به جای آن از نتیجه 
$ab'\leq a'b$
استفاده می‌کنیم که با هم معادل هستند با این خاصیت که در نامساوی دوم هر دو طرف نامساوی عدد صحیح هستند.

\متن‌سیاه{ورودی:}\\
%در خط اول ورودی، ابتدا k و سپس m داده می‌شود.\\
در خط اول ورودی به‌ترتیب $k$، $n$ و $m$ داده می‌شود، در خط دوم ورودی اعداد $a_1,a_2,...,a_k$ داده می شوند و در $n$ خط بعدی پیشنهادها به صورت 
$t_j$، $i_j$ و $b_j$ 
داده می‌شود که $t_j$ نوع پیشنهاد $j$اُم می‌باشد.

\متن‌سیاه{خروجی:}\\
%در تنها سطر خروجی، تعداد آیفون‌های مورد نیاز می‌آید. اگر در هیچ صورتی چنین کاری ممکن نبود ۱- چاپ می‌شود.\\
در تنها خط خروجی \متن‌سیاه{باقیمانده}
$\prod_{i=1}^{k}a_i$
را بر $10^9+7$ چاپ کنید (پس از اعمال $m$ پیشنهادی که آن را ماکسیمم می‌کنند).\\

\متن‌سیاه{محدودیت‌ها:}\\
$1\leq k\leq 10^5$\\
$0\leq m\leq n\leq 10^5$\\
$1\leq a_i\leq 10^4$\\
$1\leq t_j\leq 2$\\
$1\leq i_j\leq k$\\
$1\leq b_j\leq 10^4$



%$m\leq10^9$

%\textbf{\متن‌سیاه{ورودی و خروجی نمونه}\\
\متن‌سیاه{ورودی و خروجی نمونه:}\\


%% تست اول
\begin{latin}
\begin{tabularx}{\textwidth}{|X | X|}
\hline
\textbf{stdin} &
\textbf{stdout}\\
\hline
5 4 1

2 2 5 2 5

1 5 5

1 1 3

2 1 2

2 5 2&
500\\
\hline

\end{tabularx}
\end{latin}

%% تست دوم

%اگر مسئله چند قسمت داشت به این صورت عمل کنید:
%\شروع{زیرمسائل}
%\زیرمسئله[0] % مسئلۀ ۱آ
%\زیرمسئله % مسئلۀ ۱ب
%\زیرمسئله % مسئلۀ ۱پ
%\زیرمسئله % مسئلۀ ۱ت
%\پایان{زیرمسائل}

\solution{
پیشنهاد‌های نوع ۱ برای یک $a_i$ خاص را در نظر بگیرید، اگر برحسب $b_j$هایشان از بزرگ به کوچک مرتب کنیم، حتما یک پیشوند از آن (که می‌تواند تهی باشد) انتخاب شده‌ است (وگرنه می‌توان پیشنهادهای دیگری انتخاب کرد که ضرب اعداد بیشتر شود)، فرض کنید بعد از مرتب کردن به 
$b'_1\leq b'_2\leq...\leq b'_t$
رسیدیم در این صورت با توجه به چیزی که گفتیم می‌توان این پیشنهادهای نوع اول را به پیشنهادهای نوع دوم تبدیل کرد به این‌صورت که $b$ درخواست $j$-ام آن را برابر 
$\frac{a_i+\sum_{i=1}^{j}b'_i}{a_i+\sum_{i=1}^{j-1}b'_i}$
قرار می‌دهیم.\\
پس از تبدیل کردن همه‌ی پیشنهادهای نوع ۱ به نوع ۲ با استفاده از روش بالا کافیست $m$ پیشنهادی را انتخاب کنیم که $b$ بیشتری دارند (چون هر پیشنهاد نوع ۲، 
$\prod_{i=1}^{k}a_i$
را در $b$ خودش ضرب می‌کند
).


}




%%%%%%%%%%%%%%%%%%%%%%%%%%%%%%%%%%%%%%%%%%%%%%%%%%%%%
% برخی دستورات مفید برای تمرین‌ها. %
%%%%%%%%%%%%%%%%%%%%%%%%%%%%%%%%%%%%%%%%%%%%%%%%%%%%%

% فرمول‌های نمونه:
%$f(x) = \Theta(x)$
%$T(x, y) \leq \log(x) + 2^y + \binom{2n}{n}$
% {\tt code\_function}


% لیست موارد بدون شماره‌گذاری:
%\شروع{فقرات}
% \فقره مورد اول 
% \شروع{فقرات}
% \فقره مورد اول داخلی 
% \شروع{فقرات}
% \فقره اولین مورد خیلی داخلی 
% \فقره دومین مورد خیلی داخلی 
% \پایان{فقرات}
% \فقره مورد دوم داخلی 
% \پایان{فقرات}
% \فقره مورد دوم 
%\پایان{فقرات}

% لیست موارد شماره‌گذاری شده:
%\شروع{شمارش}
% \فقره یک مورد 
% \فقره یک مورد دیگر 
% \فقره باز هم یک مورد دیگر
%\پایان{شمارش}

% برای نوشتن یک فرمول چند خطی که یک شماره بخوره. برای حالت‌های دیگر به سایت زیر مراجعه کنید:
% https://www.overleaf.com/learn/latex/Aligning_equations_with_amsmath
%\begin{align} 
% \begin{split}
% (x+y)^3 &= (x+y)^2(x+y) \\
% &= (x^2+2xy+y^2)(x+y) \\
% &= (x^3+2x^2y+xy^2) + (x^2y+2xy^2+y^3) \\
% &= x^3+3x^2y+3xy^2+y^3
% \end{split} 
%\end{align}

% نوشتن ماتریس و آرایه دوبعدی
%\begin{align}
% A = 
% \begin{bmatrix}
% A_{11} & A_{21} \\
% A_{21} & A_{22}
% \end{bmatrix}
%\end{align}

% اضافه کردن تصویر با شرح. برای اضافه کردن تصویر بدون شرح، فقط دستور درج‌تصویر را استفاده کنید.
%\شروع{شکل}[h!]
%\تنظیم‌ازوسط
% \درج‌تصویر[پهنا=0.5\پهنای‌متن]{img.jpg}
%\شرح{شرح تصویر}
%\برچسب{تصویر:۱}
%\پایان{شکل}
%\پایان{زیرمسائل}

%کد زیر را ببینید.
%%\lstinputlisting{source_filename.py}
%\begin{latin}
%\begin{lstlisting}
%# Import the modules
%import sys
%import random
%
%ans = True
%
%while ans:
%question = raw_input("Ask the magic 8 ball a question: (press enter to quit) ")
%
%answers = random.randint(1,8)
%
%if question == "":
%sys.exit()
%
%elif answers == 1:
%print "It is certain"
%
%elif answers == 2:
%print "Outlook good"
%
%elif answers == 3:
%print "You may rely on it"
%
%elif answers == 4:
%print "Ask again later"
%
%elif answers == 5:
%print "Concentrate and ask again"
%
%elif answers == 6:
%print "Reply hazy, try again"
%
%elif answers == 7:
%print "My reply is no"
%
%elif answers == 8:
%print "My sources say no"
%\end{lstlisting}
%\end{latin}